% =========================================================
% Dyadic Rational Approximation
% =========================================================

\subsection{Dyadic Rational Approximation}
\label{sec:dyadic-rational-approximation}

Dyadic rationals are the rational numbers obtained by repeated halving. They are the simplest systematic approximation grid and foreshadow binary expansions, nested brackets, and Cauchy constructions.

\begin{definition}[Dyadic rational]
\label{def:dyadic-rational}
A rational number $d$ is called \emph{dyadic} if there exist $m\in\mathbb{Z}$ and $n\in\mathbb{N}$ such that
\[
d=\frac{m}{2^n}.
\]
\end{definition}

\begin{remark*}[Interpretation]
Dyadic rationals are the grid points produced by subdividing the unit interval into $2^n$ equal pieces.
\end{remark*}

\begin{tcolorbox}[colback=thmbox, colframe=thmborder, arc=2pt,
  left=6pt, right=6pt, top=4pt, bottom=4pt,
  title={\small\textbf{Theorem (Dyadic Rationals Are Rational)}},
  fonttitle=\small\bfseries]
\begin{theorem}[Dyadic Rationals Are Rational]
\label{thm:dyadic-rationals-are-rational}
Every dyadic rational is an element of $\mathbb{Q}$.
\end{theorem}
\end{tcolorbox}

\begin{remark*}[Status]
Missing proof. This follows from closure of $\mathbb{Q}$ under integer embedding and division by nonzero powers of $2$.
\end{remark*}

\begin{tcolorbox}[colback=thmbox, colframe=thmborder, arc=2pt,
  left=6pt, right=6pt, top=4pt, bottom=4pt,
  title={\small\textbf{Theorem (Dyadic Midpoint Closure)}},
  fonttitle=\small\bfseries]
\begin{theorem}[Dyadic Midpoint Closure]
\label{thm:dyadic-midpoint-closure}
If $a$ and $b$ are dyadic rationals, then $(a+b)/2$ is dyadic.
\end{theorem}
\end{tcolorbox}

\begin{remark*}[Status]
Missing proof. Put the two dyadics over a common denominator $2^N$ and then divide the sum by $2$.
\end{remark*}

\begin{tcolorbox}[colback=thmbox, colframe=thmborder, arc=2pt,
  left=6pt, right=6pt, top=4pt, bottom=4pt,
  title={\small\textbf{Theorem (Dyadic Approximation Error Bound)}},
  fonttitle=\small\bfseries]
\begin{theorem}[Dyadic Approximation Error Bound]
\label{thm:dyadic-approximation-error-bound}
For every $x\in\mathbb{Q}$ and every $n\in\mathbb{N}$, there exists $m\in\mathbb{Z}$ such that
\[
\frac{m}{2^n}\le x<\frac{m+1}{2^n}.
\]
Consequently,
\[
0\le x-\frac{m}{2^n}<\frac{1}{2^n}.
\]
\end{theorem}
\end{tcolorbox}

\begin{remark*}[Status]
Missing proof. This uses the integer floor/well-ordering step for the grid point immediately below $2^n x$.
\end{remark*}
